\section{Data Pipeline}
\label{sec:data}

Our pipeline transforms 56 textbook-extracted seed questions into 866 verified training examples through four stages: seed extraction, synthetic generation with cross-model verification, cleaning, and paraphrase augmentation. Figure~\ref{fig:pipeline} illustrates the complete flow.

\begin{figure*}[t]
\centering
\resizebox{\textwidth}{!}{%
\begin{tikzpicture}[
    node distance=0.8cm and 0.55cm,
    every node/.style={font=\small},
    box/.style={rectangle, draw, rounded corners=3pt, minimum height=0.8cm,
                align=center, inner sep=4pt},
    data/.style={box, fill=blue!10, font=\small\bfseries, minimum width=2cm},
    proc/.style={box, fill=orange!12, text width=2.2cm},
    model/.style={box, fill=green!10, text width=2.2cm},
    train/.style={box, fill=red!10, minimum width=2cm},
    arrow/.style={-{Stealth[length=5pt]}, thick},
    every label/.style={font=\scriptsize},
]

% === Row 1: Seed Extraction (horizontal) ===
\node[proc] (ocr) {Textbook OCR\\(32 chapters)};
\node[proc, right=of ocr] (extract) {Regex + LLM\\extraction};
\node[proc, right=of extract] (curate) {Manual\\curation};
\node[data, right=of curate] (seed) {56 seeds};

\draw[arrow] (ocr) -- (extract);
\draw[arrow] (extract) -- (curate);
\draw[arrow] (curate) -- (seed);

% === Row 2: Cross-model verification (horizontal, below row 1) ===
\node[model, below=1.2cm of ocr] (gen) {Claude Opus 4\\generates $(q,r,a)$};
\node[model, right=of gen] (verifier) {GPT-5.4\\solves $q$};
\node[proc, right=of verifier] (compare) {Compare\\answers (5\%)};
\node[data, right=of compare] (verified) {158 verified};

\draw[arrow] (gen) -- (verifier);
\draw[arrow] (verifier) -- (compare);
\draw[arrow] (compare) -- (verified);
% seed feeds down into generation
\draw[arrow] (seed.south) -- ++(0,-0.35) -| (gen);

% === Row 3: Cleaning (horizontal, below row 2) ===
\node[proc, below=1.2cm of verifier, text width=4.7cm] (cleanproc) {Merge, classify, split multi-part,\\deduplicate, standardize answers};
\node[data, below=1.2cm of verified] (master) {215 numeric};

\draw[arrow] (cleanproc) -- (master);
% verified feeds down into cleaning (enters from the top)
\draw[arrow] (verified.south) -- ++(0,-0.35) -| (cleanproc.north);
% seed feeds into cleaning via bayonet path: down, left past Stage 3 label, down, then right into .west
\draw[arrow, rounded corners=5pt]
    (seed.south) -- ++(0,-0.5)
    -| ([xshift=-1.8cm]cleanproc.west)
    |- (cleanproc.west);

% === Row 4: Augmentation (horizontal, below row 3) ===
\node[model, below=1.2cm of cleanproc, text width=2.2cm] (hard) {Hard generation\\Opus 4.6 + GPT-5.4};
\node[model, right=1.0cm of hard, text width=2.2cm] (para) {Paraphrase aug.\\Opus 4.6 + GPT-5.4};
\node[data, below=1.2cm of master] (v4) {866 examples};

\draw[arrow] (hard) -- node[above, font=\scriptsize] {+65} (para);
\draw[arrow] (para) -| (v4);
% master feeds down into both augmentation paths (drop well below cleanproc)
\draw[arrow] (master.south) -- ++(0,-0.7) -| (hard);
\draw[arrow] (master.south) -- ++(0,-0.7) -| (para);

% === Row 5: Training (below row 4) ===
\node[train, below=1.0cm of v4] (sft) {SFT (Qwen3-8B, QLoRA)};
\node[train, right=0.8cm of sft] (grpo) {GRPO};

\draw[arrow] (v4) -- (sft);
\draw[arrow] (sft) -- (grpo);

% === Stage labels (left side) ===
\node[font=\scriptsize\bfseries, text=black!50, left=0.2cm of ocr, anchor=east] {Stage 1};
\node[font=\scriptsize\bfseries, text=black!50, left=0.2cm of gen, anchor=east] {Stage 2};
\node[font=\scriptsize\bfseries, text=black!50, left=0.2cm of cleanproc, anchor=east] {Stage 3};
\node[font=\scriptsize\bfseries, text=black!50, left=0.2cm of hard, anchor=east] {Stage 4};
\node[font=\scriptsize\bfseries, text=black!50, left=0.2cm of sft, anchor=east] {Stage 5};

\end{tikzpicture}%
}
\caption{Data pipeline overview. Blue: dataset milestones; orange: processing; green: LLM-based generation/verification; red: training.}
\label{fig:pipeline}
\end{figure*}

\subsection{Seed Extraction}
\label{sec:data:seed}

We digitized 32 chapters from two reliability engineering textbooks---\emph{Applied Reliability} (3rd ed.) by Tobias and Trindade~\cite{tobias2012applied} and \emph{Reliability Engineering} (3rd ed.) by Elsayed~\cite{elsayed2012reliability}---using Mistral OCR. From this corpus, we extracted \textbf{56 seed QRA triplets} via regex scanning for exercise blocks, LLM-assisted validation for self-containment, and structured extraction into $(q, r, a)$ format.

\subsection{Cross-Model Verification}
\label{sec:data:selfinst}

Unlike Self-Instruct~\cite{wang2023selfinstruct}, we use a \emph{stronger} generator (Claude Opus~4) to create training data for the smaller Qwen3-8B target. Each generation prompt includes 3 randomly sampled seeds and requests 2 new problems with full chain-of-thought reasoning (temperature 0.8). Quality gates filter out short questions, placeholder answers, and near-duplicates. This yielded 158 synthetic QRA triplets.

\subsubsection{Verification Pipeline}
\label{sec:data:verification}

To ensure correctness, we replace same-model consistency with cross-model verification: (1)~Claude Opus~4 generates $(q,r,a)$ from seeds; (2)~GPT-5.4 independently solves $q$ without seeing the answer or reasoning; (3)~numerical answers are compared within 5\% tolerance, with LLM-based semantic comparison as fallback. Only matching items are retained. Because the verifier is a different model with different failure modes, agreement provides a stronger correctness signal than same-model consistency.

\subsubsection{Self-Instruct Ceiling Effect}

We observe that baseline models achieve 92--97\% accuracy on self-instruct-generated questions while scoring only 68--76\% on textbook problems, meaning same-model answer consistency~\cite{wang2023selfinstruct} is trivially satisfied. We term this the \emph{self-instruct ceiling effect}: generators create questions within their own capability envelope, avoiding edge cases they cannot solve. This motivates our paraphrase strategy.

\subsection{Cleaning and Augmentation}
\label{sec:data:augment}

The combined seed and synthetic data underwent LLM-assisted cleaning using Claude Opus~4.6: type classification, splitting of 107 multi-part questions into self-contained entries, answer standardization (stripping units, normalizing LaTeX), and deduplication by prefix matching. This yielded \textbf{215 numeric questions}. We then augment via two strategies:

\begin{itemize}
    \item \textbf{Hard generation:} 65 new difficult questions (Claude Opus~4.6, verified by GPT-5.4).
    \item \textbf{Paraphrase augmentation} (inspired by MetaMath~\cite{yu2024metamath}): Claude Opus~4.6 generates surface-level rephrases (changing wording, scenario framing) while strictly preserving all numerical values and formulas. Each rephrase is verified by GPT-5.4 for numerical and mathematical equivalence.
\end{itemize}

The contrast is striking: adding 65 hard questions yields only $+2.0\%$, while 586 paraphrases yield $+18.6\%$ ($p=0.002$), confirming that surface diversity outperforms difficulty scaling.

\subsection{Dataset Summary}
\label{sec:data:summary}

\begin{table}[t]
\centering
\small
\caption{Dataset versions. All entries are verified numeric QRA triplets.}
\label{tab:datasets}
\begin{tabular}{@{}lrc@{}}
\toprule
Version & Size & Source \\
\midrule
v1 (cleaned) & 215 & 98 textbook + 158 verified \\
v2 (+ hard)  & 280 & v1 + 65 hard generated \\
v3 (+ para.) & 501 & v2 + 221 paraphrased \\
\textbf{v4 (+ para.)} & \textbf{866} & \textbf{v2 + 586 paraphrased} \\
\bottomrule
\end{tabular}
\end{table}
